(1) 单磁矩配分函数 $z=e^{\beta\mu B}+e^{-\beta\mu B}=2\cosh x$，$x=\mu B/(k_BT)$。每个磁矩在场方向的平均投影为 $\mu\tanh x$，因此总磁化量 $M=N\mu\tanh x$，方向随 $B$ 选择。
(2) 当 $x\ll1$，$\tanh x=x+O(x^3)$，故 $M\simeq N\mu^2B/(k_BT)$，零场静磁化率 $\chi_B=N\mu^2/(k_BT)$，符合Curie反比温度定律；这是对总磁矩的响应系数，体磁化率还需除以样品体积。
(3) 内能 $U=-N\mu B\tanh x$，定场热容 $C_B=(\partial U/\partial T)_B=Nk_Bx^2\operatorname{sech}^2x\ge0$。熵 $S=Nk_B[\ln(2\cosh x)-x\tanh x]$，若固定 $B>0$ 且 $T\to0$，$x\to\infty$，熵和热容均趋零。
(4) 自旋熵仅取决于 $x=\mu B/(k_BT)$，正磁场与温度区中可逆等熵要求 $x=\mathrm{const}$，故 $B/T=\mathrm{const}$、$T_f/T_i=B_f/B_i$。形式上 $B\to0$ 可使 $T\to0$，但实际核自旋、晶格、残余相互作用与有限时间热漏限制可达到的最低温度。
